\section{Background}
%\todo[inline]{In this section, describe the background for the topic of your thesis. Generally, make use of the features of documents (and LaTeX) by providing figures, tables, referring, and other abilities to structure the text.}

After the informal introduction we will now formally introduce the various relevant technologies for the proposed thesis.

\subsection{Brief Overview of 3D Gaussian Splatting (3DGS)}
3DGS is a young AI technology in the field of 3D Reconstruction which is itself part of the wider Computer Graphics field. It was first published by Kerbl et al.\ at INRIA in France in 2023 \cite{kerbl3DGaussianSplatting2023} and has since experienced wide adoption and received critical acclaim across domains, having been cited more than 10,000 times \cite{semanticScholarKerbl,researchGateKerbl}.
3DGS solved issues previous technologies in 3D Reconstruction had, such as Neural Radiance Fields (NeRFs)~\cite{mildenhallNeRFRepresentingScenes2020}, as for example being able to render 1080p resolution scenes in real-time (\textit{Novel-View Synthesis}), making full use of modern GPGPUs, offering editability of the scene while still featuring very reasonable training times.
3DGS uses 3D Gaussians as the base primitive, a previously known data structure \cite{zwickerEWAVolumeSplatting2001,zwickerSurfaceSplatting2001}, which can be described as ellipsoid fuzzy blobs that can take any ellipsoid shape from perfect spheres to thin and long shapes or flat and round shapes.
They use a multi-dimensional normal distribution for their opacity distribution, being most opaque at the center and rapidly losing opacity towards the edges. Kerbl et al., among other achievements, made 3D Gaussians differentiable so they can be used in standard backpropagation via gradient descent.
This makes 3DGS an explicit, learnable radiance field representation without neural networks.

\subsection{3D Gaussian Splatting Algorithm}
Regular 3DGS takes 2D input images and learns the pictured 3D scene from them. The training phase can run for 2-20 minutes, depending on scene size but already good results are available within seconds after the start of training \cite{kerbl3DGaussianSplatting2023}. A large 3D scene can have millions of 3D Gaussians and hundreds of input images.
The input images are first converted into a point cloud and camera matrices (camera orientation in space etc.)\ using an algorithm called \textit{Structure from Motion} \cite{ullmanInterpretationStructureMotion1979} and the points of the point cloud become the origins for the initial 3D Gaussians before training.
It is important to know that even initialising 3D Gaussians randomly still converges, but takes longer to train.
Each forward and backward pass processes one random input image to adjust the 3D scene iteration by iteration until the loss between original input images and rendered images is small enough.
The backward pass adjusts all attributes of a 3D Gaussian (position, opacity, shape, color) purely by the difference in pixel color between rendered and original image after the previous forward pass \cite{kerbl3DGaussianSplatting2023, kerblRepo}.

I attended the AIG seminar \enquote{Artificial Intelligence - Visual Computing} with 3DGS as the topic, covering the base concept and algorithm in much more detail.

\subsection{Dynamic Scene Reconstruction}
Dynamic scene reconstruction (as in \enquote{reconstructing a dynamic scene}) in general represents reconstruction of 3D scenes that change over time, which means the input is a video feed with a certain frame rate. 

\subsubsection{Offline Dynamic Scene Reconstruction}
In the context of 3DGS as a technology, dynamic scene reconstruction is also called \enquote{4DGS} (e.g. \cite{yangDeformable3DGaussians2024, wu4DGaussianSplatting2024,yangRealtimePhotorealisticDynamic2024}). Training the initial scene is identical to regular 3DGS, but the following frames can rely on scene geometry from the previous frame, making training a lot faster as 3D Gaussians are already in the vicinity, with similar shape and color of the next frame's content.
It can still not be done in real-time, though, as the input imagery is still full 2D images and the full 3D scene is being considered for training.
The difference of the 3D scene between frames can be encoded into/learned by so-called \textit{deformation fields}, small neural networks which allow transforming 3D Gaussians over time according to the field~\cite{wu4DGaussianSplatting2024}. The other option is to encode the time component directly into the 3D Gaussian data structure, effectively transforming them into 4D Gaussians~\cite{yangRealtimePhotorealisticDynamic2024}.
This means that 4DGS scenes can only be learned from pre-recorded imagery, then encoded and then later replayed and explored with novel-view synthesis.
There is no real-time component during reconstruction\footnote{Novel-view synthesis remains real-time, as with regular 3DGS}, which is why the attribute \enquote{offline} can be assigned to this scenario.

\subsubsection{Online and Real-Time Dynamic Scene Reconstruction}
Contrary to 4DGS scenes, online dynamic scene reconstruction (also called \enquote{3DGS Live Streaming}, \enquote{Live 3DGS}, or \enquote{3DGS Live Video}) describes concurrent reconstruction and novel-view synthesis, usually after initial regular training of a base 3DGS scene.
Online dynamic scene reconstruction does not have to be real-time, but it can be, which means that real-time dynamic scene reconstruction represents the strictest class in dynamic scene reconstruction.

In both cases, a set of new input imagery is converted into the next iteration of the full 3D scene using various algorithms on the fly. There is no need for deformation fields or 4D Gaussians, even though they have been used in this context to store the dynamic scene for later.

Real-time dynamic scene reconstruction has so far not been achieved for unconstrained/\allowbreak arbitrary scenes and has seen some solutions imposing major constraints on the allowed geometry.
Existing solutions e.g.\,use AI priors (large pre-trained neural networks) to predict the next movement in the scene, limiting what can appear in the scene, or use skeletal models of humans to guide 3D Gaussians.

\subsubsection{Camera Setups for Dynamic Scene Reconstruction}
In general, dynamic scene reconstruction scenarios can be divided into single or multi-camera setups. Single cameras can be found in 3DGS solutions to SLAM (Simultaneous Localization and Mapping, mobile camera) in robotics and live 3D telepresence (fixed camera) for example. 

For the proposed thesis a fixed multi-camera setup is used in order to be able to compare consecutive frames in a meaningful way. Mobile cameras add much to the general computational load as their position has to be estimated with every frame, which is why using fixed multi-camera setups can help bound the load. This automatically limits the extent of the scene but comes naturally for e.g.\,confined indoor spaces.

Another option within this differentiation is so-called \textit{event-based 3DGS} which uses \enquote{event cameras} that only record changes in light intensity, without color. State-of-the-art solutions for SLAM or robotics in general use these kind of cameras allowing real-time 3D reconstruction of unconstrained scenes, however, at lower resolutions and without color. The proposed thesis takes a different approach, but solutions in SLAM, especially with event-based 3DGS, could prove to be valuable input, too.

Another option fascilitating dynamic reconstruction are RGB-D cameras, which record depth information alongside color values. This however works only for small scenes as larger scenes require Lidar-based sensors, which have their own restrictions and side-effects.
